\begin{solapartat}
El treball fet per la força provinent del camp elèctric serà igual a la variació de l'energia cinètica dels ions de $Xe^{+}$ \punts{0,2} $\Rightarrow$
\[ F\,d = E\,q\,d = \frac{1}{2}\,m\,v^2 \ \text{\punts{0,2}} \Rightarrow \]
\[ F = \frac{1}{2}\,\frac{mv^2}{d} = m\,a \Rightarrow a = \frac{1}{2}\,\frac{v^2}{d} = 4,5\cdot 10^{11}\ \mathrm{m/s^2} \ \text{\punts{0,2}} \]
\begin{align*}
E = \frac{1}{2}\,\frac{mv^2}{qd} \ \text{\punts{0,2}} &= \frac{1}{2}\,132\ \mathrm{u}\times\frac{1,66\cdot 10^{-27}\ \mathrm{kg}}{1u}\times\frac{(3\cdot 10^{5})^2\ \mathrm{m^2/s^2}}{1,6\cdot 10^{-19}\ \mathrm{C}\times 0,1\ \mathrm{m}}\\
&= 6,16\cdot 10^{5}\ \mathrm{N/C} \ \text{\punts{0,2}}
\end{align*}
\end{solapartat}

\begin{solapartat}
Al ser un camp elèctric constant tindrem:
\begin{align*}
&E = -\frac{\Delta V}{\Delta x} \ \text{\punts{0,3}} \Rightarrow \Delta V = -\Delta x\,E = -\frac{1}{2}\,\frac{mv^2}{q} \ \text{\punts{0,3}}\\
&\Rightarrow V_+ - V_- = \frac{1}{2}\,\frac{mv^2}{q} = 6,16\cdot 10^{4}\ \mathrm{V} \ \text{\punts{0,2}}
\end{align*}
Si la velocitat de sortida dels ions és la mateixa encara que la separació entre plaques sigui més petita, de la última expressió veiem que la diferència de potencial entre les plaques és independent de la seva separació, per tant la diferència de potencial serà la mateixa tant si $d = 10\ \mathrm{cm}$ com si és $d = 6\ \mathrm{cm}$ \punts{0,2}
\end{solapartat}
